\begin{table}[htp]
	\centering
		\begin{tabular}{ | m{7cm} | m{8cm} | }
			\hline
			\rowcolor[HTML]{5B9BD5} 
			\textbf{Requisito Funcional} & \textbf{Descripción}\\
			\hline
			
			RF-01: Procesado y Filtrado de Datos & El sistema debe cargar los datos provenientes del \textit{dataset} e identificar y descartar aquellos registros físicamente inconsistentes \\
			\hline
			RF-02: Transformación y Escalado & El sistema debe aplicar técnicas de normalización cuando sea necesario sobre las variables. Esta transformación debe ajustarse única y estrictamente sobre los conjuntos de datos de entrenamiento para prevenir el filtrado de información de los conjuntos de evaluación \\
			\hline
			RF-03: Entrenamiento de Algoritmos Estáticos & El sistema debe permitir instanciar, configurar y entrenar modelos basados en árboles de decisión, concretamente \textit{Random Forest} y \textit{XGBoost}, capaces de predecir simultáneamente variables con naturalezas distintas \\
			\hline 
			RF-04: Implementación de Redes Neuronales Multipropósito & El sistema debe ser capaz de construir arquitecturas de \textit{Deep Learning} con múltiples salidas independientes, permitiendo predecir simultáneamente variables con naturalezas distintas \\
			\hline
			RF-05: Arquitecturas Jerárquicas & El sistema debe soportar topologías de red secuenciales que respeten la causalidad física \\
			\hline
			RF-06: Búsqueda de Hiperparámetros & El sistema debe integrar módulos de optimización para explorar combinaciones de hiperparámetros que minimicen la función de pérdida seleccionada \\
			\hline
			RF-07: Validación Metodológica & El sistema debe dividir los datos empleando técnicas de retención para aislar un conjunto de validación final completamente ciego \\
			\hline
			RF-08: Cálculo de Métricas y Evaluación & El sistema debe calcular métricas de error absolutas y relativas directamente sobre las magnitudes reales para mostrar de forma objetiva y representativa la eficiencia de los modelos \\
			\hline
			RF-09: Visualización de Resultados & El sistema debe generar las gráficas pertinente para la contrastación de los resultados obtenidos por los modelos con el conjunto de validación final \\
			\hline
			
		\end{tabular} 
		
		\caption{Requisitos de información del sistema}
	\end{table}